\documentclass{article}
\usepackage{amsmath}
\usepackage{caption}
\usepackage{fontspec}
\setmainfont{WenQuanYi Zen Hei}

\begin{document}
\title{Tool Report}
\author{姓名-学号}
\maketitle

\section{Result}

本文采用的累加公式如式\eqref{eq:sum}所示，可用于计算自然数序列的求和结果。
\begin{equation}
\label{eq:sum}
\sum_{i=1}^{n} i = \frac{n(n+1)}{2}
\end{equation}

实验样本的统计数据如表\ref{tab:sample}所示，包含两组测试的核心指标。
\begin{table}[htbp]
    \centering
    \caption{实验样本数据}
    \label{tab:sample}
    \begin{tabular}{|c|c|}
        \hline
        样本组 & 数值 \\
        \hline
        Group A & 128 \\
        Group B & 256 \\
        \hline
    \end{tabular}
\end{table}

\end{document}
